% Part: first-order-logic
% Chapter: proof-systems
% Section: seqeunt-calculus

\documentclass[../../../include/open-logic-section]{subfiles}

\begin{document}

\iftag{FOL}
      {\olfileid{fol}{prf}{seq}}
      {\olfileid{pl}{prf}{seq}}

\olsection{د سېکوېنټ حساب}

د !!{derivation} ډېر نظامونه د !!{sentence}s له ترتيبونو سره کار
کوي، خو د سېکوېنټ حساب له \emph{سېکوېنټونو} سره کار کوي۔ سېکوېنټ
د دې بڼې يو عبارت دے:
\[
!A_1, \dots, !A_m \Sequent !B_1, \dots, !B_n,
\]
\textbf{د ژباړې د نسخې سمون (OLPF-001):} سرچينه د سېکوېنټ دواړو
خواوو ته يو شان وروستے اندېکس کاروي، چې دواړه لړۍ په ناپامۍ يو شان
اوږدوي؛ خو ورپسې څرګندوي چې هره لړۍ په خپلواکه توګه تشه کېدے شي۔
دلته د ښۍ لړۍ وروستے اندېکس بېل شوے، څو دواړه متناهي لړۍ خپلواکه
اوږدوالی ولري۔
يعنې دا د !!{sentence}s د لړيو يوه جوړه ده چې د سېکوېنټ په
نښه~$\Sequent$ بېله شوې ده۔ له دواړو لړيو هره يوه تشه کېدے شي۔ د
سېکوېنټ په حساب کښې !!^a{derivation} د سېکوېنټونو يوه ونه ده۔ تر
ټولو پاسني سېکوېنټونه ځانګړې بڼه لري، او ``لومړني سېکوېنټونه'' يا
``بديهي اصول'' بلل کېږي؛ هر بل سېکوېنټ د استنتاج د کومې قاعدې له
لارې سم له پاسه سېکوېنټونو څخه راځي۔ د استنتاج قاعدې يا په
سېکوېنټونو کښې !!{sentence}s سمبالوي، يعنې کيڼې يا ښۍ خوا ته ئې
ورزياتوي، ترې لرې کوي يا ئې بيا ترتيبوي؛ او يا د قاعدې په پايله کښې
يو مرکب !!{formula} معرفي کوي۔ مثلاً د $\LeftR{\land}$ قاعده اجازه
راکوي چې له $!A, \Gamma \Sequent \Delta$ څخه $!A \land !B, \Gamma
\Sequent \Delta$ استنتاج کړو، او $\RightR{\lif}$ اجازه راکوي چې د
هر $\Gamma$، $\Delta$، $!A$ او~$!B$ دپاره له $!A, \Gamma \Sequent
\Delta, !B$ څخه $\Gamma \Sequent \Delta, !A \lif !B$ استنتاج کړو۔
(په تېره بيا، $\Gamma$ او~$\Delta$ تش کېدے شي۔)

د سېکوېنټ پر حساب ولاړه $\Proves$ اړيکه داسې تعريفېږي:
$\Gamma \Proves !A$ هله او يوازې هله چې داسې کومه لړۍ $\Gamma_0$
موجوده وي چې په $\Gamma_0$ کښې هر $!A$ په~$\Gamma$ کښې وي، او داسې
!!{derivation} موجود وي چې سېکوېنټ~$\Gamma_0 \Sequent !A$ ئې په
جرړه کښې وي۔ $!A$ هله د سېکوېنټ په حساب کښې قضيه ده چې سېکوېنټ
~$\Sequent !A$ کوم !!a{derivation} ولري۔ مثلاً دا !!a{derivation}
ښيي چې $\Proves (!A \land !B) \lif !A$:
\begin{prooftree}
  \Axiom$!A \fCenter !A$
  \RightLabel{\LeftR{\land}}
  \UnaryInf$!A \land !B \fCenter !A$
  \RightLabel{\RightR{\lif}}
  \UnaryInf$\fCenter (!A \land !B) \lif !A$
\end{prooftree}

يو سټ $\Gamma$ د سېکوېنټ په حساب کښې هله ناسازګار دے چې د
$\Gamma_0 \Sequent$ کوم !!a{derivation} موجود وي، چې په کښې هر
$!A \in \Gamma_0$ په~$\Gamma$ کښې وي او د سېکوېنټ ښۍ خوا تشه وي۔
د \RightR{\Weakening} قاعدې په کارولو له ناسازګار سټ څخه هر
!!{sentence} !!{derive}d کېدے شي۔

د سېکوېنټ حساب په ۱۹۳۰مو کلونو کښې ګيرهارډ ګېنتڅن جوړ کړ۔ د منظمې
او متناظرې طرحې له امله دا د !!{derivation}s د نظريې د ودې دپاره
ډېر ګټور صوري نظام دے۔ د سېکوېنټ په حساب کښې د !!{derivation}s
موندل نسبتاً اسانه دي، خو دا !!{derivation}s ډېر ځله په سختۍ لوستل
کېږي او له ثبوتونو سره ئې تړاو کله کله په اسانۍ نۀ ښکاري۔ بيا هم،
دا د !!{derivation} نظامونو ډېره ښکلې لار ثابته شوې، او ډېر منطقونه
د سېکوېنټ حساب نظامونه لري۔

\end{document}
